\RequirePackage{newfloat}
\RequirePackage{tcolorbox}
\RequirePackage{empheq}
\providecommand{\Kernel}{\text{null}}
\providecommand{\Rfree}{R_f} 
\providecommand{\Rsf}{r_{\SDF e}}
\providecommand{\RetOp}{\mathbf{R}}
\providecommand{\R}{\mathbb{R}}
\providecommand{\N}{\mathbb{N}}
\providecommand{\Real}{\mathbb{R}}
\providecommand{\SDF}{m}
\providecommand{\Sec}{\R^N}
\providecommand{\Ret}{\mathcal S} 
\providecommand{\abs}[1]{\left| #1 \right|}
 
\providecommand{\tracemath}{\operatorname{trace}}
\providecommand{\traceop}{\operatorname{tr}}
\providecommand{\teps}{\tilde \epsilon} 
%-------------------------
% GENERAL MODIFIERS
%-------------------------
% Support both legacy \trace{A} and operator-style \trace A / \trace(A).
\ProvideDocumentCommand{\trace}{g}{%
  \IfNoValueTF{#1}{\traceop}{{\tracemath}\left( #1 \right)}%
}
\providecommand{\EST}[1]{\hat{#1}} 
\providecommand{\TIL}[1]{\tilde{#1}} 
\providecommand{\VEC}[1]{\mathbf{#1}} 
% \providecommand{\VEC}[1]{\boldsymbol{#1}} %math bold cursive 
\ifdefined\symbf
  \providecommand{\VECG}[1]{\symbf{#1}}
\else
  \providecommand{\VECG}[1]{\boldsymbol{#1}}
\fi
%-------------------------
% BOLD
%-------------------------
\providecommand{\port}{\VEC w}
\providecommand{\portw}{\VEC w}
\providecommand{\portv}{\VEC v} 
\providecommand{\portu}{\VEC u}
\providecommand{\bA}{{\VEC{A}}}
\providecommand{\bB}{{\VEC{B}}}
\providecommand{\bC}{{\VEC{C}}}
\providecommand{\bD}{{\VEC{D}}}
\providecommand{\bE}{{\VEC{E}}}
\providecommand{\bF}{{\VEC{F}}}
\providecommand{\bG}{{\VEC{G}}}
\providecommand{\bH}{{\VEC{H}}}
\providecommand{\bI}{{\VEC{I}}}
\providecommand{\bJ}{{\VEC{J}}}
\providecommand{\bL}{{\VEC{L}}}
\providecommand{\bLambda}{{\VECG{\Lambda}}}
\providecommand{\bM}{{\VEC{M}}}
\providecommand{\bN}{{\VEC{N}}}
\providecommand{\bP}{{\VEC{P}}}
\providecommand{\bQ}{{\VEC{Q}}}
\providecommand{\bR}{{\VEC{R}}}
\providecommand{\bS}{{\VEC{S}}}
\providecommand{\bU}{{\VEC{U}}}
\providecommand{\bV}{{\VEC{V}}}
\providecommand{\bW}{{\VEC{W}}}
\providecommand{\bX}{{\VEC{X}}}
\providecommand{\bY}{{\VEC{Y}}}
\providecommand{\bZ}{{\VEC{Z}}}
\providecommand{\ba}{{\VEC{a}}}
\providecommand{\bb}{{\VEC{b}}} 
\providecommand{\bc}{{\VEC{h}}}
\providecommand{\bd}{{\VEC{d}}}
\providecommand{\be}{{\VEC{e}}}
\providecommand{\bff}{{\VEC{f}}}
\providecommand{\bgg}{{\VEC{g}}}
\providecommand{\bh}{{\VEC{h}}}
\providecommand{\bone}{{\VEC{1}}}
\providecommand{\bp}{{\VEC{p}}}
\providecommand{\br}{{\VEC{r}}}
\providecommand{\bs}{{\VEC{s}}}
\providecommand{\bu}{{\VEC{u}}}
\providecommand{\bv}{{\VEC{v}}}
\providecommand{\bw}{{\VEC{w}}} 
\providecommand{\bx}{{\VEC{x}}}
\providecommand{\by}{{\VEC{y}}}
\providecommand{\bz}{{\VEC{z}}}
\providecommand{\balpha}{{\VECG{\alpha}}}
\providecommand{\bbeta}{{\VECG{\beta}}}
\providecommand{\bgamma}{{\VECG{\gamma}}}
\providecommand{\bepsilon}{{\VECG{\epsilon}}}
\providecommand{\bfeta}{{\VECG{\eta}}}
\providecommand{\blambda}{{\VECG{\lambda}}}
\providecommand{\bmu}{{\VECG{\mu}}}
\providecommand{\bnu}{{\VECG{\nu}}}
\providecommand{\bdelta}{{\VECG{\delta}}}
\providecommand{\bDelta}{{\VECG{\Delta}}}
\providecommand{\btheta}{{\VECG{\theta}}}
\providecommand{\brho}{{\VECG{\rho}}}
\providecommand{\bphi}{{\VECG{\phi}}}
\providecommand{\bpsi}{{\VECG{\psi}}}
\providecommand{\bTheta}{{\VECG{\Theta}}}
\providecommand{\bPhi}{{\VECG{\Phi}}}
\providecommand{\bPsi}{{\VECG{\Psi}}}
\providecommand{\bxi}{{\VECG{\xi}}}
\providecommand{\bSigma}{{\VECG{\Sigma}}}
\providecommand{\bOmega}{{\VECG{\Omega}}} 
\providecommand{\bzero}{{\VEC{0}}}

%-------------------------
% CAL
%-------------------------
\providecommand{\calA}{{\mathcal{A}}}
\providecommand{\calB}{{\mathcal{B}}}
\providecommand{\calD}{{\mathcal{D}}}
\providecommand{\calF}{{\mathcal{F}}}
\providecommand{\calL}{{\mathcal{L}}}
\providecommand{\calS}{{\mathcal{S}}}
\providecommand{\calX}{{\mathcal{X}}}
\providecommand{\calY}{{\mathcal{Y}}}
\providecommand{\calZ}{{\mathcal{Z}}}
\providecommand{\positives}{{\mathcal{P}}}
\providecommand{\corr}{\text{cor}}
\providecommand{\eqnperiod}{\;.}
\providecommand{\eqncomma}{\;,}
\providecommand{\variance}{\text{var}}
\providecommand{\cov}{\text{cov}}
\providecommand{\hadamard}{\circ}

%-------------------------
% TILDE AND HAT
%-------------------------
\providecommand{\tr}{\tilde r} 
\providecommand{\tu}{\tilde u} 
\providecommand{\tv}{\tilde v} 
\providecommand{\tw}{\tilde w} 
\providecommand{\tA}{\tilde A} 
\providecommand{\tB}{\tilde B} 
\providecommand{\tU}{\tilde U} 
\providecommand{\tV}{\tilde V}  
\providecommand{\tD}{\tilde D} 
\providecommand{\tLambda}{\tilde{\Lambda}} 
\providecommand{\tsigma}{\tilde{\sigma}}
\providecommand{\hr}{\hat{r}} 
\providecommand{\hf}{\hat{f}} 
\providecommand{\hg}{\hat{g}} 
\providecommand{\hatbw}{\hat{\bw}}

\providecommand{\suchthat}{\operatorname{{s.}{t.}}} 
\providecommand{\rank}{\operatorname{rank}}
\providecommand{\argmin}{\operatorname*{arg\,min}}
\providecommand{\argmax}{\operatorname*{arg\,max}}
\providecommand{\dom}{\operatorname{dom}}
\providecommand{\epi}{\operatorname{epi}}
\providecommand{\Var}{\operatorname{Var}}
\providecommand{\Cov}{\operatorname{Cov}}
\providecommand{\Bias}{\operatorname{Bias}}
\providecommand{\MSE}{\operatorname{MSE}}
\providecommand{\df}{\operatorname{df}}
\providecommand{\range}{\operatorname{range}}
\providecommand{\vecc}{\operatorname{vec}}
\providecommand{\E}{\mathbb{E}}
\providecommand{\Prob}{\mathbb{P}}
\providecommand{\Sym}{\mathbb{S}}
\providecommand{\PD}{\Sym^{p}_{++}}
\providecommand{\PSD}{\Sym^{p}_{+}}
\providecommand{\inner}[2]{\left\langle #1,\, #2 \right\rangle}
\providecommand{\nuc}[1]{\norm{#1}_{*}}
\providecommand{\half}{\tfrac{1}{2}}
\providecommand{\eps}{\varepsilon}
\providecommand{\Q}{\mathcal{Q}}
\providecommand{\Loss}{\mathcal{L}}
\providecommand{\Data}{\mathcal{D}}
\providecommand{\Om}{\Omega}
\providecommand{\Kl}{\mathrm{KL}}

%-------------------------
% NONPARAMETRIC REGRESSION
%-------------------------
% Use \providecommand so manuscripts can retain deliberate local notation.
\providecommand{\Pp}{\mathbb{P}}
\providecommand{\cF}{\mathcal{F}}
\providecommand{\cH}{\mathcal{H}}
\providecommand{\cR}{\mathcal{R}}
\providecommand{\cD}{\mathcal{D}}
\providecommand{\indic}{\mathbf{1}}
\providecommand{\lesssimc}{\mathrel{\lesssim}}
\providecommand{\practical}[1]{%
  \par\smallskip\noindent\textbf{Practical implication.} #1\par\smallskip
}
\providecommand{\warning}[1]{%
  \par\smallskip\noindent\textbf{Failure mode.} #1\par\smallskip
}

\providecommand{\st}{\text{s.t.}}
\providecommand{\etal}{\emph{et al.}}

\providecommand{\simmath}{\operatorname{sim}}
\providecommand{\cossim}{\simmath} 

\providecommand{\distmath}{\operatorname{dist}}
\providecommand{\dist}{\distmath} 
 
\providecommand{\sdvmath}{\operatorname{stdev}}
\providecommand{\sdv}{\sdvmath}  

\providecommand{\turnmath}{\operatorname{turn}}
\providecommand{\turn}{\turnmath}  

\providecommand{\negatives}{\mathcal{N}}
%\providecommand{\norm}[1]{\left|\left|#1 \right|\right|}
\providecommand{\norm}[1]{\left\lVert#1\right\rVert}
\providecommand{\proj}{\text{proj}} 
\providecommand{\sign}[1]{\text{sgn} \left(#1 \right)}
% Support both legacy \diag{d} and operator-style \diag d / \diag(d).
\ProvideDocumentCommand{\diag}{g}{%
  \IfNoValueTF{#1}{\diagmath}{\diagmath\left(#1\right)}%
}
\providecommand{\dg}[1]{\text{dg} \left(#1 \right)}
\providecommand{\diagmath}{\operatorname{diag}}
\providecommand{\optimal}{^\star}%
\providecommand{\transpose}{'}%{^\top}%{^\mathsf{T}}
\providecommand{\var}{\text{var}}
\providecommand{\inv}{^{-1}} 
\providecommand{\eqdist}{\stackrel{d}{=}}
\providecommand{\refpar}[1]{(\ref{#1})}
\providecommand{\exdiff}[1]{{\{#1\}}}
\providecommand{\similar}[2]{{\text{sim}\left( #1, #2\right)}}
\providecommand{\Image}[1]{\text{Im}\left[#1]\right)}
\providecommand{\VARIMAX}{\texttt{VARIMAX}}
\providecommand{\starsection}{$\star$}
\let\oldemptyset\emptyset
\let\emptyset\varnothing


% Additional notation from Rankbooks_v_4/newcommands.tex.  These use
% \providecommand so definitions above (and definitions supplied by a class or
% package) remain authoritative when names overlap.
\providecommand{\matdet}[1]{\left| #1 \right|}
\providecommand{\hatport}{\hat\port}
\providecommand{\tildeport}{\tilde\port}

\providecommand{\bK}{{\VEC{K}}}
\providecommand{\bT}{{\VEC{T}}}
\providecommand{\btt}{{\VEC{t}}}
\providecommand{\vone}{{\VEC{e}}}
\providecommand{\bsigma}{{\VECG{\sigma}}}
\providecommand{\btau}{{\VECG{\tau}}}
\providecommand{\bchi}{{\VECG{\chi}}}
\providecommand{\bGamma}{{\VECG{\Gamma}}}
\providecommand{\bzeta}{{\VECG{\zeta}}}

\providecommand{\hatbff}{\hat{\bff}}
\providecommand{\hatbb}{\hat{\bb}}
\providecommand{\hatbr}{\hat{\br}}
\providecommand{\hatbA}{\hat{\bA}}
\providecommand{\hatbB}{\hat{\bB}}
\providecommand{\hatbC}{\hat{\bC}}
\providecommand{\hatbD}{\hat{\bD}}
\providecommand{\hatbU}{\hat{\bU}}
\providecommand{\hatbV}{\hat{\bV}}
\providecommand{\hatbW}{\hat{\bW}}
\providecommand{\hatbX}{\hat{\bX}}
\providecommand{\hatbY}{\hat{\bY}}
\providecommand{\hatbZ}{\hat{\bZ}}
\providecommand{\hatbDelta}{\hat{\VECG \Delta}}
\providecommand{\hatbLambda}{\hat{\bLambda}}
\providecommand{\hatbSigma}{\hat{\VECG \Sigma}}
\providecommand{\hatbalpha}{\hat{\balpha}}
\providecommand{\hatbbeta}{\hat{\bbeta}}
\providecommand{\hatbgamma}{{\hat{\VECG \gamma}}}
\providecommand{\hatbdelta}{{\hat{\VECG \delta}}}
\providecommand{\hatbtheta}{{\hat{\VECG \theta}}}
\providecommand{\hatbepsilon}{\hat{\VECG \epsilon}}
\providecommand{\hatbfeta}{\hat{\VECG \eta}}
\providecommand{\hatbmu}{{\hat\bmu}}

\providecommand{\tildebalpha}{{\hat{\VECG \alpha}}}
\providecommand{\tildebbeta}{{\tilde{\VECG \beta}}}
\providecommand{\tildebepsilon}{{\tilde{\VECG \epsilon}}}
\providecommand{\tildeba}{\tilde{\ba}}
\providecommand{\tildebb}{\tilde{\bb}}
\providecommand{\tildebff}{\tilde{\bff}}
\providecommand{\tildebB}{\tilde{\bB}}
\providecommand{\tildebU}{\tilde{\bU}}
\providecommand{\tildebV}{\tilde{\bV}}
\providecommand{\tildebW}{\tilde{\bW}}
\providecommand{\tildebX}{\tilde{\bX}}
\providecommand{\tildebY}{\tilde{\bY}}
\providecommand{\tildebZ}{\tilde{\bZ}}
\providecommand{\tildebr}{\tilde{\br}}
\providecommand{\tildebs}{\tilde{\bs}}
\providecommand{\tildebu}{\tilde{\bu}}
\providecommand{\tildebv}{\tilde{\bv}}
\providecommand{\tildebw}{\tilde{\bw}}
\providecommand{\tildebx}{\tilde{\bx}}
\providecommand{\tildeby}{\tilde{\by}}
\providecommand{\tildebz}{\tilde{\bz}}

\providecommand{\COV}{{\bOmega}}
\providecommand{\hatCOV}{\hat{\VECG \Omega}}
\providecommand{\hatCOVret}{\hat{\VECG \Omega}_\br}
\providecommand{\hatCOVfac}{\hat{\VECG \Omega}_\bff}
\providecommand{\hatCOVeps}{\hat{\VECG \Omega}_\bepsilon}
\providecommand{\tildeCOVret}{\tilde{\VECG \Omega}_\br}
\providecommand{\tildeCOVfac}{\tilde{\VECG \Omega}_\bff}
\providecommand{\tildeCOVeps}{\tilde{\VECG \Omega}_\bepsilon}
\providecommand{\COVret}{\VECG{\Omega}_\br}
\providecommand{\COVfac}{\VECG{\Omega}_\bff}
\providecommand{\COVeps}{\VECG{\Omega}_\bepsilon}

\providecommand{\calQ}{{\mathcal{Q}}}
\providecommand{\calU}{{\mathcal{U}}}
\providecommand{\calV}{{\mathcal{V}}}

\providecommand{\vecmath}{\operatorname{vec}}
\providecommand{\vecmat}{\vecmath}
\providecommand{\msemath}{\operatorname{MSE}}
\providecommand{\mse}{\msemath}
\providecommand{\qlikemath}{\operatorname{QLIKE}}
\providecommand{\qlike}{\qlikemath}
\providecommand{\qlikedistmath}{\operatorname{QDIST}}
\providecommand{\qlikedist}{\qlikedistmath}
\providecommand{\operatormath}{\operatorname{op}}
\providecommand{\operator}{\operatormath}
\providecommand{\sharpemath}{\operatorname{SR}}
\providecommand{\sharpe}{\sharpemath}
\providecommand{\memph}[1]{\text{\emph{#1}}}
\providecommand{\opnorm}[1]{\norm{#1}}
\providecommand{\rademacher}{\hat R}


% Paragraph authorship bars: green Human and blue AI labels in the left margin.
% Usage: \begin{humantext} ... \end{humantext}
%        \begin{aitext}    ... \end{aitext}
% The body may contain multiple paragraphs, displayed mathematics, or lists.
% Blocks break across pages, repeating the vertical label on each segment.
% Text keeps its normal width; reserve space in the left margin for the label.
% Optional arguments accept tcolorbox keys; redefine HumanBarColor/AIBarColor
% after this file is loaded to customize the colors.
\ifdefined\gappyParagraphBarsLoaded
\else
\newcommand{\gappyParagraphBarsLoaded}{}
\tcbuselibrary{skins,breakable}
\definecolor{HumanBarColor}{HTML}{79CD76}
\definecolor{AIBarColor}{HTML}{70B2FF}
\tcbset{authorship paragraph/.style 2 args={
  enhanced,blanker,breakable,parbox=false,
  left=8pt,grow to left by=8pt,
  before skip=\medskipamount,after skip=\medskipamount,
  borderline west={1.2pt}{0pt}{#1},
  overlay={\node[rotate=90,anchor=south east,inner sep=0pt,
    font=\sffamily\scriptsize\bfseries,text=#1]
    at ([xshift=-3pt]frame.north west) {#2};}
}}
\newtcolorbox{humantext}[1][]{authorship paragraph={HumanBarColor}{Human},#1}
\newtcolorbox{aitext}[1][]{authorship paragraph={AIBarColor}{AI},#1}
\fi

% A project may supply its own styled float and theorem environments while
% still importing the shared notation above.
\ifdefined\gappyProjectCommandsLoaded
\else
% boxes
\DeclareFloatingEnvironment{takeaway}
\DeclareFloatingEnvironment{introbox}
\DeclareFloatingEnvironment{FAQ}
\DeclareFloatingEnvironment{procedure} 
\DeclareFloatingEnvironment{insight}  

\NewEnviron{yellowboxedequation}{%
  \begin{empheq}[box=\fcolorbox{yellow!50!black}{yellow!10}]{align}
  \BODY
  \end{empheq}%
}


\newenvironment{procedureenv}[1]
{
	\begin{procedure}[!htbp]
	\begin{tcolorbox}[sharp corners, colback=blue!10!white, colframe=blue!40!black]
	\caption{#1} 
	\rule{\textwidth}{0.4pt}
}
{
	\end{tcolorbox}
	\end{procedure}
}
\newenvironment{insightenv}[1]
{
	\begin{insight}[!htbp]
	\begin{tcolorbox}[sharp corners, colback=red!10!white, colframe=red!40!black]
	\caption{#1} 
	\rule{\textwidth}{0.4pt}
}
{
	\end{tcolorbox}
	\end{insight}
}
\newenvironment{faqenv}[1]
{
	\begin{FAQ}[!htbp]
	\begin{tcolorbox}[sharp corners,colback=green!10!white, colframe=green!40!black]
	\caption{\emph{#1}} 
	\rule{\textwidth}{0.4pt}
}
{
	\end{tcolorbox}
	\end{FAQ}
}
\newenvironment{introboxenv}
{
	\begin{introbox}[!htbp]
	\begin{tcolorbox}[sharp corners] 
		{\centering \textbf{The Questions}\\} 
	\rule{\textwidth}{0.4pt}
}
{
	\end{tcolorbox}
	\end{introbox}
}
\newenvironment{takeawayenv}
{
	\begin{takeaway}[!htbp]
	\begin{tcolorbox}[colframe=white!60!black, sharp corners] 
		{\centering \textbf{The Takeaways}\\} 
	\rule{\textwidth}{0.4pt}
}
{
	\end{tcolorbox}
	\end{takeaway}
}



% Number theorem environments by chapter in book-like classes and by section
% in article-like classes.
\makeatletter
\@ifundefined{c@chapter}{%
  \@ifundefined{lemma}{\newtheorem{lemma}{Lemma}[section]}{}
  \@ifundefined{theorem}{\newtheorem{theorem}{Theorem}[section]}{}
  \@ifundefined{proposition}{\newtheorem{proposition}{Proposition}[section]}{}
  \@ifundefined{assumption}{\newtheorem{assumption}{Assumption}[section]}{}
  \@ifundefined{corollary}{\newtheorem{corollary}{Corollary}[section]}{}
  \@ifundefined{definition}{\newtheorem{definition}{Definition}[section]}{}
  \@ifundefined{exercise}{\newtheorem{exercise}{Exercise}[section]}{}
  \@ifundefined{example}{\newtheorem{example}{Example}[section]}{}
}{%
  \@ifundefined{lemma}{\newtheorem{lemma}{Lemma}[chapter]}{}
  \@ifundefined{theorem}{\newtheorem{theorem}{Theorem}[chapter]}{}
  \@ifundefined{proposition}{\newtheorem{proposition}{Proposition}[chapter]}{}
  \@ifundefined{assumption}{\newtheorem{assumption}{Assumption}[chapter]}{}
  \@ifundefined{corollary}{\newtheorem{corollary}{Corollary}[chapter]}{}
  \@ifundefined{definition}{\newtheorem{definition}{Definition}[chapter]}{}
  \@ifundefined{exercise}{\newtheorem{exercise}{Exercise}[chapter]}{}
  \@ifundefined{example}{\newtheorem{example}{Example}[chapter]}{}
}
\theoremstyle{remark}
\@ifundefined{remark}{\newtheorem{remark}[theorem]{Remark}}{}
\theoremstyle{plain}
\makeatother
\fi
